\chapter{Эксперимент: как проверять архитектуру памяти}
\label{ch:experiment}

\section{Что именно проверяется}

Эксперимент в такой системе не должен отвечать на вопрос ``похожа ли Рин на живого
собеседника''. Это поздняя, интегральная и слишком мутная проверка. Сначала нужно
проверять более узкие вещи:
\begin{enumerate}
\item правильно ли retrieval находит и соединяет материал;
\item не теряется ли provenance;
\item умеет ли память ревизоваться, а не только накапливаться;
\item не порождает ли топология систематические искажения.
\end{enumerate}

\section{Три яруса стенда}

\paragraph{Ярус 1 — синтетическая топология.}
Графы с контролируемой структурой: SBM, BA, small-world, деревья, экспандерные
варианты. Здесь проверяется не содержание, а физика потока.

\paragraph{Ярус 2 — семантический граф.}
Эмбриональный subgraph из ConceptNet, ATOMIC и позднейших доращиваний. Здесь уже
проверяются fusion, bridge-узлы, влияние типов рёбер и базовая объяснимость.

\paragraph{Ярус 3 — поток жизни.}
Синтетически порождённые эпизоды, диалоги, обязательства, коррекции, ложные гипотезы
и их последующие опровержения. Здесь проверяется уже полная петля: запись, retrieval,
ответ, ревизия, сон.

\section{Четыре типа ground truth}

Чтобы тест не стал расплывчатым, нужно различать четыре разных ground truth.
\begin{enumerate}
\item \textbf{Топологический}: какие области графа должны быть связаны или труднодостижимы.
\item \textbf{Эпистемический}: какие утверждения истинны, ложны или должны остаться в зоне незнания.
\item \textbf{Происхождёнческий}: откуда именно возник узел и какие посылки должны его поддерживать.
\item \textbf{Операциональный}: что система должна была сделать, например поднять дедлайн или снять обязательство.
\end{enumerate}

\section{Новая добавка: retrieval-аудит как кортеж}

В расширенной версии вводится retrieval-аудит
\[
\mathcal{R}
=
\bigl(
R_a,\,
R_b,\,
R_c,\,
R_p
\bigr),
\]
где:
\begin{align*}
R_a &= \text{anchor recall},\\
R_b &= \text{bridge rate},\\
R_c &= \text{contradiction hit rate},\\
R_p &= \text{provenance coverage}.
\end{align*}

У этой схемы есть практическое достоинство: она разрезает плохой ответ на четыре
совершенно разных класса причин.

\section{Батарея ключевых тестов}

\begin{center}
\begin{tabular}{p{0.07\textwidth} p{0.35\textwidth} p{0.46\textwidth}}
\toprule
\# & \textbf{Гипотеза} & \textbf{Наблюдаемая метрика} \\
\midrule
1 & хабы перетягивают внимание & распределение посещений vs степень, доля ответов через hub-shortcut \\
2 & узкий мост душит контрсвидетельства & contradiction hit rate vs ширина межкластерного сечения \\
3 & fusion переоценивает слабые со-активации & precision принятой гипотезы после отложенной ревизии \\
4 & provenance деградирует под нагрузкой & доля узлов без support-skeleton после серии консолидаций \\
5 & deletion неполно очищает след & число осиротевших потомков и residual retrieval hits \\
6 & mood-lock искажает ответный контекст & сдвиг валентности retrieval-context относительно нейтральной базы \\
\bottomrule
\end{tabular}
\end{center}

\section{Отдельно про deletion tests}

Тесты удаления в таких системах особенно важны, потому что они проверяют не
декларацию, а дисциплину данных.

\begin{workedbox}[Минимальный честный deletion test]
В стенд вводится чувствительный факт $x$. Он должен:
\begin{enumerate}
\item попасть в event-log;
\item породить как минимум один выводной или индексный след;
\item затем быть удалён по команде;
\item перестать появляться в retrieval;
\item перестать поддерживать потомков;
\item оставить audit-trace самого факта удаления.
\end{enumerate}
Если хотя бы один из пунктов не выполняется, удаление нельзя считать реализованным.
\end{workedbox}

\section{Как измерять ревизуемость, а не только точность}

Память хороша не только тогда, когда она попадает в истину с первой попытки. Она
хороша тогда, когда умеет \emph{дёшево возвращаться} из ошибки. Поэтому рядом с
обычными precision/recall нужны динамические метрики:
\begin{enumerate}
\item средняя задержка до понижения ложного узла после прихода контрсвидетельства;
\item средняя глубина каскада ревизии;
\item доля потомков, переживших понижение родителя без альтернативной опоры;
\item число ложных повторных resurrect после удаления или ретракции.
\end{enumerate}

\section{Чего нельзя проверять в самом конце}

Есть соблазн сначала построить всё, а потом задать большой разговорный сценарий и
посмотреть, «ведёт ли себя память разумно». Это полезно как финальная демонстрация,
но почти бесполезно как инструмент разработки. Если ответ не удался, непонятно,
сломался ли якорь, топология, fusion, контекст-пакер или голос.

\begin{principle}[Большой сценарий завершает, а не заменяет аналитику]
Интегральный end-to-end тест нужен, но его ценность появляется только после того,
как у системы есть локальная измеримость отдельных механизмов.
\end{principle}

\section{Что должно выйти из хорошего стенда}

Хороший экспериментальный контур даёт не только зелёные/красные лампы, но и
\emph{карту слабых мест}. После серии прогонов разработчик должен уметь сказать:
\begin{enumerate}
\item память теряет якоря, но хорошо держит bridges;
\item или память находит материал, но плохо объясняет происхождение;
\item или память богата на гипотезы, но слишком медленно их ревизует;
\item или память чисто эпистемически здорова, но топологически застревает в хабах.
\end{enumerate}

Только тогда стенд действительно служит архитектуре, а не просто производит числа.
